\section{总结与延伸}

本节把 EP101 压缩成六个判断。第一，AI 应用创业的主角正在从模型公司转向应用公司和 Agent，但应用公司不能只做 API 壳。第二，产品经理底盘来自第三方视角、体验敏感和数据基建。第三，bitter lesson 要求应用公司顺着模型能力而不是用旧脚手架遮蔽模型。第四，token 消耗速度和 per token valuation 是新的产品判断工具。第五，Agent 生态会出现 AgentRank、网络效应、身份和信任问题。第六，CEO 的自我校准会影响产品路径。

\begin{importantbox}{把 EP101 放进张小珺 AI/互联网队列}
EP101 与 EP110/EP113/EP115 构成 Agent 线索：EP110 偏技术报告，EP113 偏 agentic LLM，EP115 偏 Agent 研究哲学，EP101 则是 AI 应用创业现场，讲模型能力如何通过容器、环境和网络效应转成产品。
\end{importantbox}

\subsection{概念总表}

本节把全文概念集中起来。它不是术语装饰，而是为了让读者看到 EP101 的逻辑链：产品底盘决定如何识别体验；容器和环境决定模型如何被使用；AgentRank 和网络效应决定生态如何形成；用户大于 Ego 决定组织如何持续校准。

\noindent\begin{tabularx}{\textwidth}{p{0.18\textwidth}p{0.36\textwidth}X}
\toprule
概念 & 含义 & 本期中的作用 \\
\midrule
容器/环境 & 承载模型、用户、作品和反馈的产品场 & 替代“壳”的低估说法。 \\
Vibe Coder & 用 AI coding 进行创作的新用户群 & 类似手机摄影师之于相机。 \\
AgentRank & Agent 调用和信任排序机制 & 解释未来任务分发和网络效应。 \\
OS Agent & 调度环境和垂直 Agent 的上层智能体 & 可能成为下一代入口。 \\
per token valuation & 每个 token 创造的用户价值 & 衡量智能消耗效率。 \\
用户大于 Ego & 用真实用户校准创始人和团队观点 & 防止自我膨胀和错误方向。 \\
\bottomrule
\end{tabularx}

\subsection{后续观察问题}

\begin{enumerate}
\item YouWare 能否从 HTML-to-website 的创作入口，扩展成持续创作和分发环境？
\item Vibe Coder 是短期新鲜感，还是会像手机摄影师一样形成稳定新人群？
\item AgentRank 是否会真的出现，还是被现有搜索、应用商店和模型入口吸收？
\item OS Agent 会由模型公司、操作系统公司、浏览器公司还是应用创业公司拿到？
\item Agent 身份、权限、信任和审计是否会成为新的基础设施赛道？
\item YouWare 如何在模型持续变强的情况下，把环境和经验变成可防守壁垒？
\end{enumerate}

\subsection{与 Manus、Lovart 的对照}

前面列出后续观察问题，本节把 EP101 放进 2025 上半年 Agent 应用三部曲，目的是避免把所有 Agent 创业都看成同一种形态。Manus 更像通用 Agent 叙事，强调世界不是线性外推、要做博弈中的变量；Lovart 是垂直设计 Agent，强调专业工作流、创作和设计师场景；YouWare 则把焦点放在 coding 创作的新容器和 OS Agent 的可能性。三者共同说明：Agent 应用不是同一种产品，它们在入口、环境、用户、分发和壁垒上都不同。

\noindent\begin{tabularx}{\textwidth}{p{0.17\textwidth}p{0.30\textwidth}X}
\toprule
案例 & 核心位置 & 与 YouWare 的关系 \\
\midrule
Manus & 通用 Agent，强调任务执行和变量位置 & YouWare 不直接复制通用 Agent，而是寻找创作环境。 \\
Lovart & 设计垂直 Agent，强调专业工作流 & YouWare 同样垂直，但切口是 coding 创作和发布。 \\
YouWare & 创作容器和潜在 OS Agent & 重点在新人群、作品载体、AgentRank 和网络效应。 \\
\bottomrule
\end{tabularx}

\begin{importantbox}{三部曲的共同结论}
2025 年 Agent 应用的核心，不是“谁先包装一个聊天框”，而是谁能找到新的任务环境、新的人群和新的分发机制。Manus、Lovart、YouWare 各自押的是不同层级：通用任务、专业设计、coding 创作容器。
\end{importantbox}

\subsection{拓展阅读}

\begin{itemize}
\item 对 Agent 技术报告和上下文工程感兴趣，可对照 EP110 Kimi K2 / ChatGPT Agent / Qwen3-Coder 技术报告笔记。
\item 对 Agentic LLM 和模型公司边界感兴趣，可对照 EP113 杨植麟 Kimi K2 访谈。
\item 对 Agent 研究哲学、interface 和 reward 感兴趣，可对照 EP115 姚顺雨访谈。
\item 对 AI 应用创业窗口感兴趣，可对照 Lovart / 陈冕和 Manus / 肖弘访谈。
\end{itemize}

\end{document}
